\section{Introduction}
A spectral analyzer shares its host's audio budget even when it produces no
audio. Its average cost may be small, yet a straightforward implementation
concentrates window preparation, the transform, and display-data processing
into whichever sample call completes a frame. VCV Rack makes the consequence
concrete: modules process one sample at a time, and the engine's worker
threads synchronize between samples~\cite{rackengine}, so one long call can
hold up every worker. Whether reducing that concentration improves the
behavior of a whole block, however, also depends on the engine's own
synchronization overhead and on the rest of the patch.

Several execution strategies are available. A batch FFT on the engine thread
minimizes state and can exploit optimized kernels. A worker thread moves the
computation elsewhere, at the cost of a capture, queueing, and ownership
policy. When an application controls the relative phases of several
analyzers, staggering them reduces coincident frame work. Cooperative
execution offers a further choice: retain an unfinished analysis and advance
it during successive sample calls. It is most relevant when a producer must
maintain an exact-hop spectral history independently of a display that polls
intermittently. For a latest-frame visualizer that tolerates dropped analyses,
computing on the UI thread is also reasonable; that is a different delivery
contract.

Making only the FFT resumable is not enough. Frame-sized packing,
reconstruction, and output passes remain at its boundaries, and restarting
each transform as soon as the previous one finishes ties the frame cadence to
rounding in the butterfly quota. We instead schedule one dependency-ordered
sequence that covers the whole analyzer, and we drive publication from a frame
clock that is independent of when the arithmetic finishes. This report
contributes:
\begin{itemize}
\item an explicit execution contract for work credits, exact-hop publication,
retention of the selected input, cancellation, and ownership of complete
spectra (Section~\ref{sec:one-hop});
\item matched batch/distributed controls and native-layout batch/hybrid
alternatives that separate work placement from FFT-kernel choice
(Section~\ref{sec:evaluation});
\item a same-host study of native completion horizons, shipped module
defaults, and paced execution through Rack's actual multithreaded engine
(Section~\ref{sec:comparison-results}).
\end{itemize}

The evidence supports a qualified conclusion. Distributing a fixed scalar
implementation substantially reduces its typical per-hop peak. Native hybrid
pipelines can combine small peaks with much lower total cost than
butterfly-level scheduling. Inside the complete host, however,
synchronization and release lateness can dominate, and a smaller local burst
does not by itself mean fewer missed releases. The practical contribution is
therefore a measured way to choose execution granularity and publication age,
with analyzer-level and host-level outcomes kept distinct.

The implementation serves two VCV Rack modules: the Fourier spectrum analyzer,
which processes four independent SIMD channels, and the scalar Spectre
spectrogram. Section~\ref{sec:related} positions the work and
Section~\ref{sec:signal-model} fixes notation. Section~\ref{sec:legacy-motivation}
explains why scheduling only the FFT falls short, and Section~\ref{sec:one-hop}
states the contract. Sections~\ref{sec:evaluation}
and~\ref{sec:comparison-results} describe the study and its results, and
Section~\ref{sec:discussion} discusses design guidance and limits. The
appendices collect derivations, an analysis of the original scheduler, and
implementation details.
